% concatenated from On_the_Theory_of_Bias_Tuning_in_Event_Cameras_5.tex
%% LyX 2.3.7 created this file.  For more info, see http://www.lyx.org/.
%% Do not edit unless you really know what you are doing.
\documentclass[english]{article}
\usepackage[T1]{fontenc}
\usepackage[latin9]{inputenc}
\usepackage{float}
\usepackage{textcomp}
\usepackage{amsmath}
\usepackage{amsthm}
\usepackage{amssymb}
\usepackage{graphicx}
\usepackage{wasysym}

\makeatletter
%%%%%%%%%%%%%%%%%%%%%%%%%%%%%% Textclass specific LaTeX commands.
\numberwithin{equation}{section}
\numberwithin{figure}{section}
\theoremstyle{plain}
\newtheorem{thm}{\protect\theoremname}
\theoremstyle{remark}
\newtheorem{rem}[thm]{\protect\remarkname}
\theoremstyle{plain}
\newtheorem{cor}[thm]{\protect\corollaryname}
\theoremstyle{definition}
\newtheorem{defn}[thm]{\protect\definitionname}
\theoremstyle{plain}
\newtheorem{lem}[thm]{\protect\lemmaname}

\makeatother

\usepackage{babel}
\providecommand{\corollaryname}{Corollary}
\providecommand{\definitionname}{Definition}
\providecommand{\lemmaname}{Lemma}
\providecommand{\remarkname}{Remark}
\providecommand{\theoremname}{Theorem}

\begin{document}
\title{On the Theory of Bias Tuning in Event Cameras}
\author{David El-Chai Ben-Ezra, Daniel Brisk, Adar Tal}
\maketitle
\begin{abstract}
This paper lays the foundation of a theory for bias tuning in neuromorphic
cameras, a novel sensing technology also known as \textquotedblleft event
cameras''. We begin by formulating the high-level effect of the sensitivity
biases on the camera\textquoteright s event rate in mathematical terms.
We then show that, as a corollary of the Poincar�--Miranda theorem,
the commonly used tuning principles of rate budgeting and polarity
balancing lead to a unique configuration of the sensitivity biases.
As a corollary, we show how by adopting these principles, the multi-variable
bias-tuning problem reduces to a two-parameter problem that can be
resolved experimentally.
\end{abstract}
Key words: neuromorphic camera, event camera, biases tuning, optimization
problem, Poincar�--Miranda theorem.\\
\\
AMS subject classification: 49J21, 93C35, 93B52, 93C65.

\section{Introduction}

An event camera is a bio-inspired sensor that represents a paradigm
shift in visual data acquisition. Instead of capturing full frames
at fixed intervals, each pixel operates independently and triggers
only when it detects a brightness change exceeding a threshold. The
output is thus a stream of asynchronous, timestamped events. Event
cameras provide microsecond temporal resolution, low latency, high
dynamic range (over 120 dB), and low power consumption, making them
potentially suited for tasks such as detection, tracking, motion segmentation,
optical flow estimation, 3D reconstruction, pose estimation, and SLAM
\cite{key-1,key-2}, especially in high-speed scenarios.

Because event cameras differ fundamentally from frame-based cameras,
significant effort has been devoted to developing asynchronous, event-driven
processing methods (see \cite{key-3}-\cite{key-24}, \cite{key-25}-\cite{key-28}).
Beyond the output format, they also differ in how data generation
is controlled. Conventional cameras rely on intuitive parameters such
as frame rate and exposure time, guided by well-established principles
like the Nyquist--Shannon theorem. Event cameras, however, do not
use such parameters. Instead, they rely on several biases that control
event generation. Two biases determine pixel sensitivity to illumination
changes, another sets an internal low-pass filter (considered as controlling
the bandwidth \cite{key-1}), a further bias governs an internal high-pass
filter (often overlooked in the literature), and a final bias defines
the pixel refractory period (see $\varoint$\ref{sec:The-biases-of}).
Because these biases directly affect the event stream, proper tuning
is essential for achieving optimal performance. Yet their effects
are less intuitive, and systematic tuning methods remain relatively
immature despite existing studies \cite{key-11}-\cite{key-23}.

A key challenge in developing tuning methods is the lack of direct
information linking bias settings to the generated data. Unlike frame-based
cameras, where photon counts can be inferred from gray levels and
exposure time, an event does not directly indicate the magnitude of
the triggering illumination change. This is due to the logarithmic
sensing mechanism, which depends on the pixel\textquoteright s initial
illumination and varies across pixels. The influence of the internal
filters further complicates this relationship, and the absence of
frames limits intuition. Therefore, there is a need for tuning methods
based on quantitative, user-accessible measurements.

In this paper we examine two high-level principles proposed for tuning
event camera sensitivity: polarity balancing \cite{key-16} and event-rate
limiting \cite{key-13,key-20}. Since equal sensitivity to increases
and decreases in illumination cannot be set directly, polarity balancing
is typically achieved by matching positive and negative event rates.
Hence, a key advantage of these principles is that they rely on parameters
easily accessible to the user. Event-rate limiting is conceptually
similar to controlling the frame rate in conventional cameras, while
stepping forward and maintaining a constant event rate resembles automatic
gain control (AGC), balancing information capture with preventing
sensor saturation.

The paper aims to establish a foundation for event-camera bias tuning,
starting from the principles of polarity balancing and event-rate
control, and extending them toward a practical heuristic for the multi-variable
tuning problem. We first describe schematically how the biases influence
the event rate and express these effects with simple mathematical
interpretations (see $\varoint$\ref{sec:Sensitivity-bias-modeling}).
The description is partly based on our experience with the Prophesee
EVK4\nobreakdash-HD, but the underlying properties also follow conceptually
from how event cameras operate. We note that in this paper, we do
not focus on validating these properties but rather on applying their
implications. We say that an event camera that satisfies these properties
in a specific scene has ``standard behavior''. 

Based on this interpretation and the Poincar�--Miranda theorem \cite{key-29},
we derive the following fundamental theorem ($\varoint$\ref{sec:The-Existence-and}):
\begin{thm}
\label{thm:main}Given a scene, an event camera with standard behavior,
and a choice of the filter biases, let $R_{P}(p,n)$ and $R_{N}(p,n)$
be the positive and negative event-rates as function of the sensitivity
biases $p$ and $n$, respectively. Then, there exists a constant
$K$ such that for every target event-rate $k\geq K$, there exists
a unique pair $(p_{0},n_{0})$ such that 
\begin{equation}
R_{P}(p_{0},n_{0})=R_{N}(p_{0},n_{0})=k.\label{eq:main}
\end{equation}
\end{thm}

\begin{rem}
We note that the theorem only proves the existence of balanced configuration,
and does not relate to the algorithmic way to converge to it. 
\end{rem}

As a corollary, we obtain the following observation:
\begin{cor}
\label{cor:Given-a-specific}Under the assumptions of Theorem \ref{thm:main},
given a target event-rate $k$, for every choice of the filter biases
$h,l$, there exists up to one unique pair $p_{0}(h,l),n_{0}(h,l)$
such that Equation \ref{eq:main} holds.
\end{cor}

\begin{rem}
In practice, there exist scene conditions and filter-bias configurations
for which the minimal achievable event rate exceeds the desired target
event rate. In such cases, Equation \ref{eq:main} admits no solution.
However, under sufficiently aggressive filter-bias configurations,
the lower bound $K$ from Theorem \ref{thm:main} is typically close
to zero, indicating that the desired target event rate can still be
achieved.
\end{rem}

As a result, under appropriate assumptions about how the target signal
depends on the sensitivity biases, the following corollary can be
derived:
\begin{cor}
\label{cor:Let--be}Let $Sig(h,l,p,n)$ be a standard target signal,
as function of the filter-biases $h,l$ and sensitivity biases $p,l$.
Then, for optimizing the signal under the principles of polarity balancing
and event-rate limiting, it is enough to examine 
\[
\tilde{Sig}(h,l)=Sig(h,l,p_{0}(h,l),n_{0}(h,l))
\]
 over the filter-biases $h,l$ plane, where $p_{0}(h,l),n_{0}(h,l)$
are the functions from Corollary \ref{cor:Given-a-specific}. .
\end{cor}

This corollary implies that, under applying the two high-level principles
above, the multi-variable problem of event-camera bias tuning can
be reduced to a two-dimensional problem that can be explored experimentally.
To illustrate this heuristic, Section $\varoint$\ref{sec:demonstration}
analyzes the two-dimensional case for a periodic light signal driven
by the electrical grid. The analysis shows that the optimal filter
bias values differ significantly from the default settings recommended
by Prophesee. In particular, the proposed heuristic can increase the
detected signal by a factor of 2--4 compared with the default configuration.

Since the goal of this paper is to explain the heuristic rather than
fully validate it, the demonstration does not cover all possible signal
and background conditions. For specific detection tasks, the study
must be adapted accordingly. Broader validation across different scenarios
will be addressed in future work focusing on particular applications.
\begin{rem}
Throughout this paper, the term \textquotedblleft event camera\textquotedblright{}
refers to models that include the tuning parameters discussed above,
excluding devices without these degrees of freedom. For clarity and
concreteness, we focus our discussion on the Prophesee EVK4-HD model.
However, the same heuristic can be adapted to other relevant models
with appropriate adjustments.
\end{rem}


\section{Event Camera Biases\label{sec:The-biases-of}}

Contrary to frame-based cameras, the output of an event camera is
not a sequence of synchronous intensity frames. Instead, it consists
of an asynchronous stream of \textquotedblleft events.\textquotedblright{}

Each pixel maintains a reference intensity level and continuously
monitors changes relative to this reference. When the change in intensity
exceeds a predefined threshold, the pixel updates its reference level
and generates an event. Each event encodes the timestamp (with microsecond
resolution), the pixel coordinates, and the polarity, indicating whether
the change was positive or negative. The camera output is thus a time-ordered
list of such events.

Prophesee event cameras include five parameters that control event
generation, each of which can be tuned upward or downward. Increasing
a parameter in one direction typically leads to more events from the
signal, but also increases background activity and noise; tuning it
in the opposite direction has the reverse effect. We now describe
these parameters.

Two parameters control the thresholds for event generation: one for
positive events and one for negative events. In Prophesee\textquoteright s
interface, these are denoted \textquotedblleft bias\_diff\_on\textquotedblright{}
and \textquotedblleft bias\_diff\_off,\textquotedblright{} which we
refer to as $p$ (positive) and $n$ (negative), respectively. These
parameters can be tuned independently; lowering them reduces the triggering
thresholds, increases sensitivity, and results in more signal events
as well as increased noise.

Two additional parameters control internal low-pass (LPF) and high-pass
(HPF) filters within each pixel. After converting incoming light into
an electrical signal, the pixel passes this signal through both filters
before evaluating whether an event should be generated. These parameters
are labeled \textquotedblleft bias\_fo\textquotedblright{} (LPF) and
\textquotedblleft bias\_hpf\textquotedblright{} (HPF), which we denote
by $l$ and $h$, respectively.

In general, increasing these parameters shifts the frequency response
toward higher frequencies. Increasing $l$ allows more high-frequency
components to pass, leading to more signal events but also higher
background activity. Increasing $h$, on the other hand, suppresses
low-frequency components, reducing both signal events and noise. The
precise relationship between these biases and the effective frequency
response, however, is not well characterized.

The final parameter controls the refractory period of each pixel---the
minimum time between consecutive events. Its main purpose is to suppress
spurious activity caused by internal noise and to limit events from
\textquotedblleft hot pixels,\textquotedblright{} which generate abnormally
high activity. This mechanism is not intrinsic to event-camera operation
and ideally would be unnecessary with perfect pixel stability. In
practice, for the EVK4-HD model, this effect is quite low: the default
refractory period is relatively short (on the order of 10--100 microseconds
\cite{key-17}), and the contribution of hot pixels is negligible.
Therefore, we do not consider this parameter in this work.

We are thus left with four degrees of freedom: two controlling sensitivity
$(p,n)$ and two governing the internal filtering $(l,h)$.

\section{Sensitivity Bias Modeling\label{sec:Sensitivity-bias-modeling}}

We start with some definitions. The term ``scene'' will denote the
environment within the camera\textquoteright s field of view whose
spatial structure, motion, and illumination produce the brightness
variations that trigger events in the sensor. 

The average numbers of positive and negative events generated by scene-induced
brightness variations and internal pixel noise are referred to as
the ``positive event-rate'' and ``negative event-rate'', denoted
by $R_{P}$ and $R_{N}$, respectively. These rates are influenced
by factors such as object motion, flickering light sources, sensor
noise, and small edge displacements caused by air turbulence. Under
this notation, polarity balancing corresponds to the condition $R_{P}=R_{N}$,
while event-rate limiting corresponds to $R_{P},R_{N}\leq k$, for
some predefined bound $k$. The choice of the event-rate bound $k$
should be determined by the user based on various factors, such as
the application requirements, camera stability, computational capacity,
and real-time processing constraints. A detailed characterization
of this bound is beyond the scope of this paper.

Since the biases govern event generation, they also determine the
event rates. Accordingly, we treat $R_{P}(l,h,p,n)$ and $R_{N}(l,h,p,n)$
as functions of the biases. Although, in practice, these parameters
are discrete, we model them as continuous for the purpose of analysis
and assume that the event rates are differentiable with respect to
the biases.

We further assume that increasing sensitivity leads to more information
from the scene, and thus $R_{P}$ and $R_{N}$ are monotonically decreasing
functions of $p$ and $n$. That is, for any choice of $l,h,p,n$ 

\begin{align*}
\frac{\partial R_{P}}{\partial p},\frac{\partial R_{N}}{\partial n} & <0\qquad\mathrm{and}\qquad\frac{\partial R_{P}}{\partial n},\frac{\partial R_{N}}{\partial p}\leq0.
\end{align*}

At first glance, it is not obvious that changing $p$ affects $R_{N}$,
or vice versa. However, lowering a sensitivity threshold has two effects
that increase the event rate. The first is that the camera becomes
responsive to smaller illumination changes. This effect primarily
influences the event rate of the corresponding polarity. In particular,
as a sensitivity bias approaches its minimal value, the corresponding
threshold approaches zero, leading to an unbounded increase in the
associated event rate. Formally, for any $l,h,p,n$

\[
\underset{n\to n_{min}}{\lim}R_{N}(l,h,p,n)=\underset{p\to p_{min}}{\lim}BER_{P}(l,h,p,n)=\infty
\]
where $p_{min}$ and $n_{min}$ denote the lower bounds of the parameters
$p$ and $n$, respectively.

The second effect is that lowering a sensitivity threshold allows
the camera to capture a larger fraction of the actual illumination
changes. Unlike the first effect, this influences the event rates
of both polarities. Consequently, reducing even a single sensitivity
bias might increase both $R_{P}$ and $R_{N}$. However, the combination
of both effects is not symmetric. As a rule of thumb, adjusting the
sensitivity of one polarity has a stronger relative impact on the
corresponding event rate than on the opposite one. This can be expressed
as

\begin{align*}
\frac{\left(\frac{\partial R_{P}}{\partial p}\right)}{R_{P}} & =\frac{\partial\ln(R_{P})}{\partial p}<\frac{\left(\frac{\partial R_{N}}{\partial p}\right)}{R_{N}}=\frac{\partial\ln(R_{N})}{\partial p}\leq0\\
\frac{\left(\frac{\partial R_{N}}{\partial n}\right)}{R_{N}} & =\frac{\partial\ln(R_{N})}{\partial n}<\frac{\left(\frac{\partial R_{P}}{\partial n}\right)}{R_{P}}=\frac{\partial\ln(R_{P})}{\partial n}\leq0.
\end{align*}

To strengthen the robustness of the analysis, we adopt a weaker assumption.
Specifically, we assume the existence of a monotonic function $\alpha:\mathbb{R}^{+}\to\mathbb{R}$
such that

\begin{align*}
\frac{\partial(\alpha\circ R_{P})}{\partial p} & <\frac{\partial(\alpha\circ R_{N})}{\partial p}\leq0\\
\frac{\partial(\alpha\circ R_{N})}{\partial n} & <\frac{\partial(\alpha\circ R_{P})}{\partial n}\leq0.
\end{align*}

We summarize the above properties in the following definition.
\begin{defn}
Given a scene, an event camera is said to exhibit ``standard behavior''
if the relationship between the sensitivity biases $p,n$ and the
positive and negative event rates $R_{P}$ and $R_{N}$ satisfies
the following properties for any choice of $l,h,p,n$:
\end{defn}

\begin{enumerate}
\item $\underset{n\to n_{min}}{\lim}R_{N}=\underset{p\to p_{min}}{\lim}R_{P}=\infty$.
\item There exists a monotonic function $\alpha:\mathbb{R}^{+}\to\mathbb{R}$
such that 
\begin{align*}
\frac{\partial(\alpha\circ R_{P})}{\partial p} & <\frac{\partial(\alpha\circ R_{N})}{\partial p}\leq0\\
\frac{\partial(\alpha\circ R_{N})}{\partial n} & <\frac{\partial(\alpha\circ R_{P})}{\partial n}\leq0.
\end{align*}
\end{enumerate}
Our final assumption concerns the dependence of the target signal
on the sensitivity biases. In this context, the ``signal'' refers
to the component of the measured data generated by the phenomenon
of interest, which carries the information to be detected or analyzed.
While the precise definition and measurement of the signal may vary
across applications, it can, in all cases, be modeled as a function
of the biases, denoted by $Sig(l,h,p,n)$. 

We assume that the information captured from the signal increases
with the sensitivity of the camera. Accordingly, we say that a signal
is ``standard'' if, for any $l,h,p,n$

\begin{equation}
\frac{\partial Sig}{\partial p},\frac{\partial Sig}{\partial n}<0.\label{eq:sig}
\end{equation}


\section{The Existence and Uniqueness Theorem\label{sec:The-Existence-and}}

Denote the minimal and maximal values of $p$ and $n$ by $p_{min},p_{max},n_{min},n_{max}$.
Given a choice of the filter biases $l,h$, let us denote the following
functions over the region $(-p_{max},-p_{min})\times(-n_{max},-n_{min})$:
\begin{align*}
f(x,y) & =R_{N}(l,h,-x,-y)\\
g(x,y) & =R_{P}(l,h,-x,-y).
\end{align*}
Then, Theorem \ref{thm:main} (and Corollary \ref{cor:Given-a-specific})
follow from the following lemma.
\begin{lem}
\label{lem:main}Let $f(x,y)$ and $g(x,y)$ be real-valued differentiable
functions, defined over the rectangles $[a,b)\times[c,d]$ and $[a,b]\times[c,d)$,
correspondingly. Assume also that for every $(x_{0},y_{0})$ in the
regions we have:

\begin{align*}
1. & \underset{x\to b}{\lim}f(x,y_{0})=\underset{y\to d}{\lim}g(x_{0},y)=\infty\\
2. & \frac{\partial(\alpha\circ f)}{\partial x}(x_{0},y_{0})>\frac{\partial(\alpha\circ g)}{\partial x}(x_{0},y_{0})\geq0\\
3. & \frac{\partial(\alpha\circ g)}{\partial y}(x_{0},y_{0})>\frac{\partial(\alpha\circ f)}{\partial y}(x_{0},y_{0})\geq0
\end{align*}
for some monotonic $\alpha:\mathbb{R}^{+}\to\mathbb{R}$. Then, there
exist two constants $0\leq K_{1},K_{2}$ such that for every $K_{1}\leq k_{1}$
and $K_{2}\leq k_{2}$ there exists a unique solution to the equation
system
\[
f(x,y)=k_{1},\qquad g(x,y)=k_{2}.
\]
\end{lem}

\begin{proof}
As we assume that $f$ and $g$ are monotonic, we have:
\begin{align*}
K_{1} & :=f(a,d)=\max\{f(a,y)\,|\,y\in[c,d]\}\\
K_{2} & :=g(b,c)=\max\{g(x,c)\,|\,x\in[a,b]\}.
\end{align*}

Now, let $K_{1}\leq k_{1}$ and $K_{2}\leq k_{2}$. Then, by Properties
1-3 there exists $b'\in[a,b)$ and $d'\in[c,d)$ such that
\begin{align*}
\min\{f(b',y)\,|\,y\in[c,d]\}= & f(b',c)>K_{1}\\
\min\{g(x,d')\,|\,x\in[a,b]\}= & g(a,d')>K_{2}
\end{align*}

Thus, by applying Poincar�--Miranda theorem \cite{key-29} on the
rectangle $[a,b']\times[c,d']$, we obtain that there exists a solution
to the equation system
\[
f(x,y)=k_{1},\qquad g(x,y)=k_{2}.
\]

For the uniqueness, notice that using the chain rule, Properties 2
and 3 yield that the of Jacobian of $H:\mathbb{R}^{2}\to\mathbb{R}^{2}$
defined by $H(x,y)=(f(x,y),g(x,y))$ is always positive:
\begin{align*}
J_{H}(x,y) & =\left|\begin{array}{cc}
\frac{\partial f}{\partial x} & \frac{\partial f}{\partial y}\\
\frac{\partial g}{\partial x} & \frac{\partial g}{\partial y}
\end{array}\right|=\frac{\partial f}{\partial x}\cdot\frac{\partial g}{\partial y}-\frac{\partial f}{\partial y}\cdot\frac{\partial g}{\partial x}\\
 & =\frac{1}{\frac{\partial\alpha}{\partial f}\cdot\frac{\partial\alpha}{\partial g}}\left(\frac{\partial(\alpha\circ f)}{\partial x}\cdot\frac{\partial(\alpha\circ g)}{\partial y}-\frac{\partial(\alpha\circ f)}{\partial y}\cdot\frac{\partial(\alpha\circ g)}{\partial x}\right)>0.
\end{align*}
Hence, $h$ is one-to-one, yielding the uniqueness of the solution,
as required. 
\end{proof}
The next lemma shows that, as one should intuitively expect, when
positive and negative event rates are balanced, increasing the target
event rate corresponds to increasing the sensitivity in both polarities.
However, somewhat unexpectedly, as will be remarked later, this relationship
does not necessarily hold when unbalanced sensitivity is required.
\begin{lem}
\label{lem:2}Let $f$ and $g$ be as in Lemma \ref{lem:main}. Then,
both coordinates of the unique solution of the system
\[
f(x,y)=g(x,y)=k
\]
grow as functions of $k$. In other words, $x(k)$ and $y(k)$ are
monotonic.
\end{lem}

\begin{proof}
Let us denote $f_{\alpha}=\alpha\circ f$ and $g_{\alpha}=\alpha\circ g$,
and write the equations
\[
f_{\alpha}(x(k),y(k))=g_{\alpha}(x(k),y(k))=\alpha(k)
\]
where $x(k)$ and $y(k)$ are the unique values given by Lemma \ref{lem:main}
for each $k$. Then, by the chain rule
\[
\begin{array}{c}
\frac{\partial\alpha}{\partial k}=\frac{\partial f_{\alpha}}{\partial k}=\frac{\partial f_{\alpha}}{\partial x}\cdot\frac{\partial x}{\partial k}+\frac{\partial f_{\alpha}}{\partial y}\cdot\frac{\partial y}{\partial k}\\
\frac{\partial\alpha}{\partial k}=\frac{\partial g_{\alpha}}{\partial k}=\frac{\partial g_{\alpha}}{\partial x}\cdot\frac{\partial x}{\partial k}+\frac{\partial g_{\alpha}}{\partial y}\cdot\frac{\partial y}{\partial k}
\end{array}\Rightarrow\left(\begin{array}{cc}
\frac{\partial f_{\alpha}}{\partial x} & \frac{\partial f_{\alpha}}{\partial y}\\
\frac{\partial g_{\alpha}}{\partial x} & \frac{\partial g_{\alpha}}{\partial y}
\end{array}\right)\cdot\left(\begin{array}{c}
\frac{\partial x}{\partial k}\\
\frac{\partial y}{\partial k}
\end{array}\right)=\left(\begin{array}{c}
\frac{\partial\alpha}{\partial k}\\
\frac{\partial\alpha}{\partial k}
\end{array}\right)
\]
yielding that
\[
\left(\begin{array}{c}
\frac{\partial x}{\partial k}\\
\frac{\partial y}{\partial k}
\end{array}\right)=\frac{1}{\frac{\partial f_{\alpha}}{\partial x}\cdot\frac{\partial g_{\alpha}}{\partial y}-\frac{\partial f_{\alpha}}{\partial y}\cdot\frac{\partial g_{\alpha}}{\partial x}}\cdot\left(\begin{array}{cc}
\frac{\partial g_{\alpha}}{\partial y} & -\frac{\partial f_{\alpha}}{\partial y}\\
-\frac{\partial g_{\alpha}}{\partial x} & \frac{\partial f_{\alpha}}{\partial x}
\end{array}\right)\cdot\left(\begin{array}{c}
\frac{\partial\alpha}{\partial k}\\
\frac{\partial\alpha}{\partial k}
\end{array}\right)
\]
and therefore, by the assumptions $\frac{\partial g_{\alpha}}{\partial y}>\frac{\partial f_{\alpha}}{\partial y}\geq0$
and $\frac{\partial f_{\alpha}}{\partial x}>\frac{\partial g_{\alpha}}{\partial x}\geq0$
we have
\begin{align*}
\frac{\partial x}{\partial k} & =\frac{\frac{\partial\alpha}{\partial k}\cdot\left(\frac{\partial g_{\alpha}}{\partial y}-\frac{\partial f_{\alpha}}{\partial y}\right)}{\frac{\partial f_{\alpha}}{\partial x}\cdot\frac{\partial g_{\alpha}}{\partial y}-\frac{\partial f_{\alpha}}{\partial y}\cdot\frac{\partial g_{\alpha}}{\partial x}}>0\\
\frac{\partial y}{\partial k} & =\frac{\frac{\partial\alpha}{\partial k}\cdot\left(\frac{\partial f_{\alpha}}{\partial x}-\frac{\partial g_{\alpha}}{\partial x}\right)}{\frac{\partial f_{\alpha}}{\partial x}\cdot\frac{\partial g_{\alpha}}{\partial y}-\frac{\partial f_{\alpha}}{\partial y}\cdot\frac{\partial g_{\alpha}}{\partial x}}>0
\end{align*}
as required.
\end{proof}
\begin{rem}
One should notice that by following the above argument, it can be
shown that Lemma \ref{lem:2} is not necessarily right under imbalanced
requirement such as $f(x,y)=k$ and $g(x,y)=\gamma\cdot k$ where
$0<\gamma\neq1$. In other words, if a user intentionally seeks an
imbalanced configuration between positive and negative event rates,
achieving a higher total event rate may require the unintuitive adjustment
of one of the sensitivity biases in the opposite direction.
\end{rem}

Lemma \ref{lem:2} implies that any standard signal will grow as function
of the total event-rate, under the principle of polarity balancing,
and standard behavior of the camera:
\begin{cor}
Given a scene, event camera with standard behavior, a standard signal
$Sig(l,h,p,n)$, and filter biases $l,h$, let $p(k)$ and $n(k)$
be the unique solution of the equation $R_{P}(p,n)=R_{N}(p,n)=k$.
Then, $Sig$ is monotonically increasing as function of $k$. 
\end{cor}

\begin{proof}
By Lemma \ref{lem:2}, the chain rule, and Equation \ref{eq:sig},
we have
\[
\frac{\partial Sig(p(k),n(k))}{\partial c}=\frac{\partial Sig}{\partial p}\cdot\frac{\partial p}{\partial k}+\frac{\partial Sig}{\partial n}\cdot\frac{\partial n}{\partial k}>0.
\]
\end{proof}
\begin{cor}
Given a scene, a standard behavior of the camera, a standard signal
$Sig(l,h,p,n)$, filter biases $l,h$, and threshold event rate $k$,
the signal reaches its highest value under the constraints
\begin{align*}
R_{P} & =R_{N}\leq k
\end{align*}
when $p(k),n(k)$ are the unique solutions of the equation $R_{P}=R_{N}=k$. 
\end{cor}

The meaning of the last corollary is that if we balance between bounding
the event rate of the camera, and having as much information from
our signal, then under the constraint of having balanced positive
and negative event-rates, the optimal signal will be reached under
the largest accepted event-rate. In particular, we obtain Corollary
\ref{cor:Let--be}, that tells us that in order to optimize the signal
under the two principles of polarity balancing and bounded event-rate
$k$, it is enough to examine the signal under the $l,p$ plane where
$p(l,h)$ and $n(l,h)$ are those biases that solve the system of
equations
\[
R_{P}(l,h,p(l,h),p(l,h))=R_{N}(l,h,p(l,h),p(l,h))=k
\]
i.e. it is enough to evaluate
\[
\tilde{Sig}(l,h)=Sig(l,h,p(l,h),n(l,h)).
\]


\section{Determining Filter Biases: A Demonstration\label{sec:demonstration} }

In this section, we present an experimental illustration of the behavior
of the functions $p(l,h)$, $n(l,h)$ and $\tilde{Sig}(l,h)$. 

For the signal, we used an incandescent lamp powered by a 50 Hz sinusoidal
electrical grid, producing a 100 Hz periodic signal of the form $|\sin(x)|$.
The signal was measured using a Prophesee EVK4-HD placed a few meters
from the lamp. To simulate a weak signal, the lamp was covered, leaving
only a small aperture for light emission. This setup produced a signal
consisting of a few dozen events per period. The laboratory environment
was kept dark to minimize background activity.

Next, we sampled the $(l,h)$ plane, corresponding to \textquotedblleft bias\_fo\textquotedblright{}
and \textquotedblleft bias\_hpf,\textquotedblright{} respectively,
in steps of 10 units. For each sample, we applied an automatic feedback
control algorithm to balance the positive and negative event rates
while maintaining a total event rate of approximately 100k events
per second (equivalent to positive and negative event-rates of $k$
= 50k event/sec each). From our experience for such scenes, such target
event-rate is higher then the lower bound $K$ required in Theorem
\ref{thm:main} for any choice of the filter-biases, and hence convergence
of the algorithm to the desired target event rate is guaranteed. 

In Figure \ref{Fig1}, the values of \textquotedblleft bias\_diff\_on\textquotedblright{}
and \textquotedblleft bias\_diff\_off\textquotedblright , namely $p(l,h)$
and $n(l,h)$, are shown for each sample. In general, these values
decrease as \textquotedblleft bias\_hpf\textquotedblright{} increases
and \textquotedblleft bias\_fo\textquotedblright{} decreases. This
indicates that when the filters are more aggressive, the background
noise is reduced, allowing the algorithm to increase the camera\textquoteright s
sensitivity by lowering the sensitivity biases while maintaining a
constant  event rate.

For each sample, after tuning the sensitivity biases, we recorded
the signal for 0.5 seconds, corresponding to 50 periods. We then counted
the number of positive and negative events within a radius of 5 pixels
around the signal center and divided by 50. Accordingly, we define
$\tilde{Sig}(l,h)$ as the average number of events per period in
this region---once for positive events and once for negative events---generated
by the periodic signal of the covered incandescent lamp.

Figure \ref{Fig2} shows the variation of $\tilde{Sig}(l,h)$ as a
function of the filter biases. It can be observed that $\tilde{Sig}$
attains its maximum values for 80$\leq$bias\_hpf$\leq$90 and \textminus 10$\leq$bias\_fo$\leq$0,
in both cases. In particular, when $\tilde{Sig}$ denotes the number
of positive events per period, its optimal value is 2--4 times higher
than the values obtained near the default settings of \textquotedblleft bias\_hpf\textquotedblright{}
and \textquotedblleft bias\_fo\textquotedblright{} (both equal to
0 in the Prophesee EVK4-HD). In the case of negative events, the optimal
value is slightly lower, reaching approximately 1.5--2 times the
default value. This difference between positive and negative responses
highlights the importance of defining $\tilde{Sig}$ appropriately
for the desired task. Additionally, the results indicate that the
signal is significantly more sensitive to \textquotedblleft bias\_hpf\textquotedblright{}
than to \textquotedblleft bias\_fo,\textquotedblright{} a somewhat
surprising observation given the limited attention paid to \textquotedblleft bias\_hpf\textquotedblright{}
in the literature.

\begin{figure}[!ph]
\includegraphics[width=5in,height=7in]{figure-on-off}\caption{Values of the sensitivity biases $p$ = bias\_diff\_on and $n$ =
bias\_diff\_off as functions of the filter biases $l$ = bias\_fo
and $h$ = bias\_hpf. For each sampled $(l,h)$, the sensitivity biases
were tuned to maintain a polarity balanced event rate of approximately
100k events/sec.\label{Fig1}}
\end{figure}

\begin{figure}[H]
\includegraphics[width=5in,height=7in]{figure-main}\caption{Number of positive and negative events per period of the signal generated
by the incandescent lamp, as a function of the filter biases.}
\label{Fig2}
\end{figure}

\begin{rem}
The tuning algorithm iterates over the biases in strips, fixing \textquotedblleft bias\_fo\textquotedblright{}
while varying \textquotedblleft bias\_hpf\textquotedblright{} in steps
of 10 units, from its minimal value up to bias\_hpf = 150. As a result,
there is a significant difference between the background event rate
at the last sample of one strip and that at the first sample of the
next strip prior to retuning. This discrepancy may cause the algorithm
to converge improperly at the beginning of each strip and explains
the observed valleys at lower values of \textquotedblleft bias\_hpf\textquotedblright . 
\end{rem}


\section{Conclusions}

In this paper, we developed a theoretical framework for bias tuning
in event cameras, motivated by the fact that the relationship between
the sensitivity biases and the generated events is often vague, making
direct tuning difficult. To overcome this, we proposed a solution
based on quantities accessible to the user, such as the event rate,
and formalized the principles of polarity balancing and event-rate
control. This approach allows the inherently multi-variable bias tuning
problem to be reduced to a structured, two-parameter problem, making
it both analyzable and experimentally tractable. A key theoretical
insight is that the existence of a balanced bias configuration follows
from the Poincar�--Miranda theorem, demonstrating that bias tuning
is grounded in rigorous mathematical principles rather than heuristics.
The notions of ``standard behavior'' for the camera and ``standard
signal'' provide a formal framework to describe how both background
activity and the target signal depend on the sensitivity biases. Experimental
results with a periodic light source illustrate that the framework
effectively guides tuning, significantly increasing the signal while
controlling background events, and highlight the important role of
the high-pass filter bias, which is often underexplored in the literature.
Overall, this work bridges theory and practice, providing a solid
mathematical basis for bias tuning in event cameras. Future work will
extend this framework to systematic validation across diverse applications.

\section*{Declaration of generative AI and AI-assisted technologies in the
manuscript preparation process}

During the preparation of this work the authors used ChatGPT for improving
language and readability. After using this tool/service, the authors
reviewed and edited the content as needed and take full responsibility
for the content of the published article.
\begin{thebibliography}{10}
\bibitem{key-1}G. Gallego, T. Delbruck, G. Orchard, C. Bartolozzi,
B. Taba, A. Censi, S. Leutenegger, A. J. Davison, J. Conradt, K. Daniilidis
and D. Scaramuzza, \textquotedbl Event-based Vision: A Survey,\textquotedbl{}
IEEE Transactions on Pattern Analysis and Machine Intelligence, 44(1),
154-180, 2020. 

\bibitem{key-2}M. H. Tayarani-Najaran and M. Schmuker, \textquotedbl Event-Based
Sensing and Signal Processing in the Visual, Auditory, and Olfactory
Domain: A Review,\textquotedbl{} Frontiers in Neural Circuits, 15,
2021.

\bibitem[3]{key-3}Z. Ni, C. Pacoret, R. Benosman, S. Ieng, S. R�GNIER,
``Asynchronous event-based high speed vision for microparticle tracking,''
Journal of Microscopy, 245(3), 236-244, 2012.

\bibitem[4]{key-4}X. Lagorce, C. Meyer, S. H. Ieng, D. Filliat, R.
Benosman, ``Asynchronous event-based multikernel algorithm for high-speed
visual features tracking'' IEEE Transactions on Neural Networks and
Learning Systems, 26(8), 1710-1720, 2014.

\bibitem[5]{key-5}I. Alzugaray, M. Chli, ACE: An efficient asynchronous
corner tracker for event cameras, IEEE 2018 International Conference
on 3D Vision (3DV), 653-661, 2018.

\bibitem[6]{key-6}I. Alzugaray, M. Chli, ``Asynchronous corner detection
and tracking for event cameras in real time,'' IEEE Robotics and
Automation Letters, 3(4), 3177-3184, 2018.

\bibitem[7]{key-7}R. Li, D. Shi, Y. Zhang, K. Li, R. Li, ``Fa-harris:
A fast and asynchronous corner detector for event cameras,'' 2019
IEEE/RSJ International Conference on Intelligent Robots and Systems
(IROS), 6223-6229, 2019.

\bibitem[8]{key-8}H. Chen, Q. Wu, Y. Liang, X. Gao, H. Wang, ``Asynchronous
tracking-by-detection on adaptive time surfaces for event-based object
tracking,'' Proceedings of the 27th ACM International Conference
on Multimedia, 473-481, 2019.

\bibitem[9]{key-9}Z. Wang, T. Molloy, P. van Goor, R. Mahony, ``Event
Blob Tracking: An Asynchronous Real-Time Algorithm,'' arXiv:2307.10593,
2023.

\bibitem[10]{key-10}D. Gehrig, D. Scaramuza, ``Low-latency automotive
vision with event cameras,'' Nature 629(8014), 1034-1040, 2024.

\bibitem[11]{key-24}D.E.C. Ben-Ezra, R. Arad, A. Padowicz, I. Tugendhaft,
``Probabilistic Approach for Detection of High-Frequency Periodic
Signals using an Event Camera,'' New Mathematics and Natural Computation,
1-12, 2024.

\bibitem[12]{key-11}P. Lichtsteiner, C. Posch, T. Delbruck, ``A
128$\times$128 120 dB 15$\mu s$ latency asynchronous temporal contrast
vision sensor,'' IEEE Journal of Solid-State Circuits, 43(2), 566-576,
2008.

\bibitem[13]{key-12}R. Gra�a, T. Delbruck, ``Unraveling the paradox
of intensity-dependent DVS pixel noise,'' International Image Sensor
Workshop (IISW), arXiv:2109.08640, 2021.

\bibitem[14]{key-13}T. Delbruck, R. Gra�a, M. Paluch, . ``Feedback
control of event cameras,'' IEEE/CVF Conference on Computer Vision
and Pattern Recognition (CVPR), 1324-1332, 2021.

\bibitem[15]{key-14}B. McReynolds, R. Gra�a, T. Delbruck, ``Experimental
methods to predict dynamic vision sensor event camera performance'',
Optical Engineering, 61(7), 074103, 2022.

\bibitem[16]{key-15}M. S. Dilmaghani, W. Shariff, C. Ryan, J. Lemley,
P. Corcoran, ``Control and evaluation of event cameras output sharpness
via bias,'' Fifteenth International Conference on Machine Vision
(ICMV 2022), SPIE ,Vol. 12701, 455-462, 2023.

\bibitem[17]{key-16}R. Gra�a, B. McReynolds, T. Delbruck, ``Shining
light on the DVS pixel: A tutorial and discussion about biasing and
optimization'', IEEE/CVF Conference on Computer Vision and Pattern
Recognition (CVPR), 4045-4053, 2023.

\bibitem[18]{key-17}P. N. McMahon-Crabtree, L. Kulesza, B. J. McReynolds,
D. S. O'Keefe, A. Puttur, D. Maestas, C. P. Morath, M. G. McHarg,
``Event-based camera refractory period characterization and initial
clock drift evaluation'', Unconventional Imaging, Sensing, and Adaptive
Optics, SPIE, Vol. 12693, 253-273, 2023.

\bibitem[19]{key-18}R. Gra�a, B. McReynolds, T. Delbruck, , ``Optimal
biasing and physical limits of DVS event noise,'' International Image
Sensor Workshop (IISW), arXiv:2304.04019, 2023.

\bibitem[20]{key-19}M.S. Dilmaghani, W. Shariff, M. A. Farooq, J.
Lemley, P. Corcoran, ``Optimization of Event Camera Bias Settings
for a Neuromorphic Driver Monitoring System'' IEEE Access, 2024.

\bibitem[21]{key-20}G. B. Nair, M. Milford, T. Fischer, ``Enhancing
Visual Place Recognition via Fast and Slow Adaptive Biasing in Event
Cameras'', arXiv preprint arXiv:2403.16425, 2024.

\bibitem[22]{key-21}V. Vishnevskiy, G. Burman, S. Kozerke, D. P.
Moeys, ``Optimal OnTheFly Feedback Control of Event Sensors'', arXiv:2408.12976,
2024.

\bibitem[23]{key-22}E. V. Votyakov, A. Artusi, ``Quantifying Noise
of Dynamic Vision Sensor'', arXiv:2404.01948, 2024.

\bibitem[24]{key-23}M.S. Dilmaghani, W. Shariff, C. Ryan, J. Lemley,
and P. Corcoran, ``Autobiasing Event Cameras'', arXiv:2411.00729.

\bibitem[25]{key-25}A. Mondal, J. H. Giraldo, T. Bouwmans, A. S.
Chowdhury, \textquotedbl Moving Object Detection for Event-based
Vision using Graph Spectral Clustering\textquotedbl , Proceedings
of the IEEE/CVF International Conference on Computer Vision, 876-884.
2021.

\bibitem[26]{key-26}S. Schaefer, D. Gehrig, D. Scaramuzza, \textquotedbl AEGNN:
Asynchronous Event-based Graph Neural Networks\textquotedbl , Proceedings
of the IEEE/CVF Conference on Computer Vision and Pattern Recognition,
12371-12381. 2022.

\bibitem[27]{key-27}K. Jeziorek, A. Pinna, T. Kryjak, \textquotedbl Memory-Efficient
Graph Convolutional Networks for Object Classification and Detection
with Event Cameras\textquotedbl , Signal Processing: Algorithms,
Architectures, Arrangements, and Applications (SPA), 160-165, IEEE,
2023.

\bibitem[28]{key-28}L. Fan, Y. Li, H. Shen, J. Li, D. Hu, \textquotedbl From
Dense to Sparse: Low-Latency and Speed-Robust Event-Based Object Detection\textquotedbl ,
IEEE Transactions on Intelligent Vehicles, 2024.

\bibitem[29]{key-29}M. N. Vrahatis, \textquotedbl A short proof
and a generalization of Miranda\textquoteright s existence theorem\textquotedbl ,
Proceedings of the American Mathematical Society 107(3), 701-703 ,1989.
\\
\\
David El-Chai Ben-Ezra, Remote Sensing Department, Soreq NRC, Yavne,
Israel 81800, dbenezra@mail.huji.ac.il\\
\\
Daniel Brisk, Remote Sensing Department, Soreq NRC, Yavne, Israel
81800, danielbr@soreq.gov.il\\
\\
Adar Tal, Remote Sensing Department, Soreq NRC, Yavne, Israel 81800,
adar.tal1@mail.huji.ac.il

\end{thebibliography}

\end{document}
